\subsection{The remaining integral ambiguity under the geometric identities}

\begin{statement}{B.38}{Global geometric identities leave eighty-one formal classes, and the commutator and local identities leave three}{Proved}
Put $M=A(-1)=\operatorname{Lie}_3(V)(-1)$ and retain the
actual arithmetic source $\Gamma$ and tangent of B.37.
Let $\mathcal A$ be the $54$ global Fermat automorphisms.
Fix a rational action class and the cohomological path
terms in all its $\mathcal A$-identities. The group of
homogeneous integral cohomology differences satisfying
these data is
\[
 \mathcal T_{\rm glob}
   =\bigl\{h\in\operatorname{Tor}H^1(\Gamma,M):
         \operatorname{res}_K h\text{ is }\mathcal A\text{-invariant}\bigr\}
       \simeq(\mathbb Z/3)^4.                               \tag{B.38.1}
\]
It has $81$ elements.

There are two additional necessary conditions for the
actual commutator extension. Define
\[
 H_0=\mathbb Z_3u+\mathbb Z_3v\subset V,
 \qquad\mathcal J=[H_0,\Lambda^2V](-1)\subset M.
                                                               \tag{B.38.2}
\]
This is a primitive rank-$90$ coefficient module, and
the actual class $\eta$ is in the image of
$H^1(\Gamma,\mathcal J)$ integrally.
The subgroup of (B.38.1) satisfying this commutator
condition is $(\mathbb Z/3)^2$.
The subgroup of (B.38.1) satisfying the local quaternion
identities at two is also $(\mathbb Z/3)^2$.
Their intersection is exactly
\[
                           \mathcal T_{\rm ref}\simeq\mathbb Z/3.
                                                               \tag{B.38.3}
\]
Thus fixing the rational Ross component, the actual
global path identities, the integral commutator condition,
and the local quaternion identities leaves precisely
three formal integral cohomology classes once one such
class is fixed. This is not a claim that all three are
realized by Galois actions of the curve.

The remaining nonzero difference is not an inner path
class. Its localization at two is, in the thirteen
coordinates of B.36,
\[
             (0,0,0,0,0;\ 0,0,0,0,0,0;\ 54,0),             \tag{B.38.4}
\]
which generates the subgroup adding $27$ to the twelfth
$\mathbb Z/81$ coordinate.
It vanishes under local coefficient reduction to
$M/3^rM$ for $r\le3$, and is detected at $r=4$ by an
explicit local coefficient quotient $M\twoheadrightarrow\mathbb Z/81$.
This specifies the finite coefficient test. B.40 computes
the actual geometric extension and evaluates that test:
the selected twelfth coordinate is $57$ modulo $81$,
which distinguishes the actual refinement from the other two.

The prescribed inherited order-$3^{30}$ quotient in
prompt B.4 kills all of $\mathcal J$. Consequently it
cannot distinguish these three formal classes, even
after the preceding constraints are imposed.
The remaining $\mathbb Z/3$ is distinct both from
LM6 E.39's $(\mathbb F_3)^3$ abelian-path kernel and
from the much larger ambient action torsion of B.21.
\end{statement}
\begin{proof}
\emph{Rational cosets describe exactly the homogeneous torsion.}
At a good Frobenius the eigenvalues on $M\otimes\mathbb Q_3$
all have positive weight. Thus
$(M\otimes\mathbb Q_3)^{G_K}=0$.
The coefficient sequence gives an injective connecting map
\[
 \delta:(M\otimes\mathbb Q_3/M)^{\Gamma}
                  \longrightarrow H^1(\Gamma,M),
 \qquad \delta(y)(g)=(g-I)y.                                 \tag{B.38.5}
\]
Its image is the kernel of rationalization, hence the
torsion subgroup. It is finite because the matrix at
one good Frobenius minus one is rationally invertible.
Equivalently, the explicit finite-generation calculation
of B.21 gives the same conclusion.

Every global geometric automorphism is $G_K$-equivariant
on the graded coefficient lattice. Naturality of (B.38.5)
and the vanishing of rational invariants show that
$\operatorname{res}_K\delta(y)$ is invariant under such
an automorphism $a$ if and only if $(a-I)y\in M$.
Therefore (B.38.1) is exactly the finite rational-coset
group
\[
                   (M\otimes\mathbb Q_3/M)^{\Gamma,\mathcal A}.
                                                               \tag{B.38.6}
\]
This identifies the homogeneous ambiguity of fixed
cohomological path identities. It does not replace
the actual inhomogeneous path terms by zero.

\emph{The global integral matrices and their exact coset calculation.}
Work first in the original CM basis of B.24, with
\[
 R=\begin{pmatrix}0&-1\\1&-1\end{pmatrix},\qquad
 C_E=\begin{pmatrix}1&-1\\0&-1\end{pmatrix}.
\]
For a compact matrix $L$ and a boundary matrix $P$, put
\[
                 B(L,P)=
       \begin{pmatrix}L&0\\(UL-PU)/3&P\end{pmatrix}.          \tag{B.38.7}
\]
The four geometric generators can be taken as two
translations by $A_0,A_2$, the CM automorphism, and
the elliptic involution. Their compact matrices are
$I,I,R,-I$; their boundary matrices are the actual
permutations of $E[3]$ modulo the diagonal, retaining
the column $(-1,\ldots,-1)^t$ whenever a cusp maps to $O$.
Formula (B.38.7) is integral for all four.
Their products give exactly $54$ distinct integral
matrices, and each preserves $H_0$ and $w$.

The arithmetic Frobenii have degree-one matrices
$B(I+3R,7I_8)$ and $B(I+3R^2,7I_8)$.
If $P_c$ swaps $A_1,A_2$, swaps $B_1,B_2$, and swaps
$C_1,C_2$, complex conjugation has degree-one matrix
$B(C_E,-P_c)$.
On $M$ the arithmetic matrices are respectively
$7^{-1}\operatorname{Lie}_3(B(I+3R,7I_8))$,
its conjugate, and $-\operatorname{Lie}_3(B(C_E,-P_c))$.
The geometric actions have no Tate factor.
All the $K$-arithmetic matrices commute with the $54$
geometric matrices, and complex conjugation permutes
the latter and exchanges the two Frobenii.

For any chosen set of these actions, stack the matrices
$g-I$ vertically. If its Smith diagonal is $3^{d_j}$,
then the invariant rational-coset group is
$\bigoplus_j\mathbb Z/3^{d_j}$.
Indeed, from an exact identity
$U_0 B V_0=\operatorname{diag}(3^{d_j})$ one obtains
the coset generators $(V_0)_j/3^{d_j}$.
There is full column rank because the seven Frobenius
minus one is rationally invertible.
The exact calculations give
\[
 \begin{array}{c|l|c}
 \text{conditions}&\text{Smith valuations}&\text{order}\\ \hline
 \Gamma&0^{237},1^{59},2^{27},3^7&3^{134}\\
 G_K,\mathcal A&0^{325},1^4,2^1&3^6\\
 \Gamma,\mathcal A&0^{326},1^4&3^4.
 \end{array}                                                  \tag{B.38.8}
\]
The first row independently reproduces B.21's ambient
torsion. The last row proves (B.38.1).
The program \path{c_global_action_ambiguity.py} uses
the four translation/CM generators and a vertical
Smith calculation with column transport.
The independent program \path{b_global_action_ambiguity.py}
constructs all $54$ matrices directly from the Fermat
plane's diagonal cube roots and coordinate permutations,
checks a different four-generator presentation, and
uses transposed rectangular Smith reduction.
Both retain complete invariant-coset generators and
give (B.38.8).

\emph{The commutator condition is integral at three.}
The primitive pair $u,v$ can be completed to a basis
of $V$, for example (B.37.6). Choose lifts of that basis
in the free pro-three group $N$.
Their Frattini images form a basis, so they freely
generate $N$: the resulting surjective endomorphism
of a finitely generated free pro-three group is an
isomorphism. Use the logarithms of these generators
in the rational class-three Lie algebra.

Only two specific degrees require integrality.
Class-two BCH uses the coefficient $1/2$, which is
integral in $\mathbb Z_3$, so the degree-two action
term $B_g:V\to\Lambda^2V$ is integral.
In $\gamma_2N/\gamma_4N$, the logarithm of a group
commutator differs from the pure bracket $[u,v]$
only by a cubic term with coefficient $1/2$.
Every cubic Lie bracket is already the logarithm
of an integral triple commutator. Subtract that
central correction to obtain an integral lift
$\omega=[u,v]$ of $w$.
This does not assert an integral class-three BCH
correspondence for the entire group; no division by
three has been made.
Formula (B.37.11) now applies with integral $B_g$.
As $A_gH_0=H_0$, it proves
\[
                  \eta(g)\in\mathcal J
                      \quad\hbox{for every }g.              \tag{B.38.9}
\]

In a basis beginning $u,v$, the map
$H_0\otimes\Lambda^2V\to\operatorname{Lie}_3V$
given by the bracket has image of rank
$2+32+56=90$.
The first two summands are the whole multidegrees
$H_0^3$ and $H_0^2(V/H_0)$.
In the last multidegree, each generator
$[u,[e_r,e_s]]$ or $[v,[e_r,e_s]]$, $r>s\ge2$,
is a difference of two distinct Hall columns with
coefficient one. The pairs of columns are disjoint
as $u$ or $v$ and $r,s$ vary.
These observations give a unit minor in each
multidegree, proving that $\mathcal J$ is primitive.
Both refinement programs also construct all $90$
columns and verify their independence already
modulo three.
The quotient $M/\mathcal J$ is free and has no
rational Galois invariants by its positive weights.
Consequently
$H^1(\Gamma,\mathcal J)\to H^1(\Gamma,M)$ is injective.
For a torsion class represented by $y$ in (B.38.5),
membership in this image is equivalent to
\[
                    y\in\mathcal J\otimes\mathbb Q_3+M.
                                                               \tag{B.38.10}
\]
This follows by applying the same rational-coset
injection to the free quotient $M/\mathcal J$.

\emph{The exact intersections, including the local comparison.}
For a concrete coordinate certificate, use the four
Smith numerators $z_0,z_1,z_2,z_3$ stored by
\path{b_global_action_ambiguity.py}; thus an element
of (B.38.1) is represented by
$y=\sum a_i z_i/3$, $a_i\in\mathbb F_3$.
Every $z_i$ modulo three is supported entirely in
$\operatorname{Lie}_3(W)$, where $W$ is the original
boundary lattice. The local strict compact conjugacy
of B.35 fixes $W$ pointwise, so fixes these particular
cosets. No unknown global/local compact change of
basis enters the following quaternion calculation.

Eliminating the $90$ primitive columns in (B.38.10)
gives the two equations
\[
                        a_1=a_3=0.                           \tag{B.38.11}
\]
For a torsion cocycle $\delta(y)$, its local value
at four is $(F_4-I)y$.
Every local quaternion automorphism $q$ commutes
with $F_4$. Rational invertibility of $F_4-I$ gives
\[
 \operatorname{res}_{(2)}\delta(y)\text{ is }q\text{-invariant}
        \quad\Longleftrightarrow\quad(q-I)y\in M.             \tag{B.38.12}
\]
Testing all eight actual quaternion matrices on the
four pure-boundary numerators gives exactly
\[
                    a_0+a_2=0,\qquad a_1+2a_3=0.             \tag{B.38.13}
\]
Each of (B.38.11) and (B.38.13) leaves a plane over
$\mathbb F_3$, hence nine classes.
Their common line is generated by $z=z_2-z_0$ modulo
three, proving (B.38.3).
This remaining numerator also has a compact geometric
description. For a cusp $P$ let $\bar e_P$ be its
boundary class in $W/3W$, with $\sum_P\bar e_P=0$.
In the free restricted Lie algebra over $\mathbb F_3$,
put
\[
 \bar z=-\sum_{\substack{L\subset E[3]\\L\text{ an affine line}}}
                \left(\sum_{P\in L}\bar e_P\right)^{[3]}.
                                                              \tag{B.38.13a}
\]
There are twelve lines. The displayed sum belongs to
the ordinary degree-three Lie subspace. Indeed in
characteristic three
$(r+s)^3-r^3-s^3=[r,[r,s]]+[s,[s,r]]$,
so cubing is additive modulo that subspace, and the
sum of the twelve line sums is $4\sum_P\bar e_P=0$.
Expanding the cubes into the $512$ words on the eight
boundary generators gives exactly the saved Hall
numerator $z$. In particular (B.38.13a), with any
integral Hall lift divided by three, specifies the
surviving cohomology generator without a choice of
Smith basis. The finite affine geometry also makes
its geometric invariance transparent.

There is an integral four-bracket lift inside
$\mathcal J$. For $(a,b)$ in
$\{(1,0),(0,1),(1,1),(1,2)\}$, set
$V_{a,b,j}=\sum_{ax(P)+by(P)=j}e_P$, where the cusp
coordinates are in $E[3]=\mathbb F_3^2$.
Put
\[
 Z=-\sum_{(a,b)}
             [a v-b u,[V_{a,b,0},V_{a,b,1}]]\in\mathcal J.
                                                              \tag{B.38.13b}
\]
For one parallel class its three boundary sums satisfy
$V_0+V_1+V_2=0$. Their cubes sum to
$[V_1-V_0,[V_0,V_1]]$ modulo three.
The columns of $U$ give
$\bar u=\sum_P y(P)\bar e_P$ and
$\bar v=-\sum_P x(P)\bar e_P$, so
$V_1-V_0=a\bar v-b\bar u$.
This proves $Z\bmod3=\bar z$ directly.
The surviving global cocycle can therefore be specified
explicitly as $g\mapsto(gZ-Z)/3$.
The independent $512$-word verification and the full
integral Hall lift are recorded by
\path{b_ambiguity_line_formula.py} and its data file.

These intersection computations were implemented in
two different integral bases. The program
\path{c_action_ambiguity_refinement.py} uses the
adapted Hall basis and two quaternion generators;
\path{b_action_ambiguity_refinement.py} constructs
$[u,[e_i,e_j]]$ and $[v,[e_i,e_j]]$ directly in the
original Hall basis and tests all eight quaternion
matrices. Both independently enumerate all $81$
coefficient vectors, finding $9$, $9$, and $3$
survivors respectively. Their different coset bases
give the same surviving line and its same oriented
localization (B.38.4).

The reduction of $z$ is not in
$\operatorname{ad}_w(V)\bmod3$, as exact elimination
against its ten primitive columns verifies.
Since the quotient has no rational invariants,
the corresponding cohomology class is not in the
inner-path image of $H^1(\Gamma,V)$ either.
On the pure boundary cubic module Frobenius four
acts by $4^3/4=16$. Choose a pure-boundary integral
Hall lift $z$ of the numerator, distinct from the
four-bracket lift $Z$ in (B.38.13b).
Hence the surviving class is
represented locally by
\[
                        (F_4-I)(z/3)=5z.                    \tag{B.38.14}
\]
The exact comparison with B.36 gives $6h_9$, namely
(B.38.4). Localization is injective on this global
torsion group: if $(F_4-I)y=(F_4-I)m$ for an integral
$m$, rational invertibility forces $y=m$.

For the final refinement it is enough to fix the
rational Ross component, rather than every rational
coordinate in advance. Indeed a difference satisfying
the same global path identities is invariant under
$\mathcal A$ in rational cohomology.
If it also lies in $H^1(\Gamma,\mathcal J)$, B.37's
multidegree argument says that its average lies
entirely in $D(-1)$. Fixing that Ross component
therefore makes the whole rational difference zero.
This justifies the formulation following (B.38.3).

\emph{The next finite coefficient test, and the inherited test's kernel.}
Let $\Pi$ be the integral quaternion projector on $M$,
$E$ its invariant basis, and $I_E$ its recorded set of
$41$ invertible minor rows. Let $Q_2$ be the row
transport in B.36, and let $r_{12}$ select the twelfth
positive Smith coordinate. The integral row
\[
 \lambda_{81}
     =r_{12}Q_2(E_{I_E})^{-1}\Pi_{I_E}\pmod{81}              \tag{B.38.15}
\]
is primitive and satisfies
$\lambda_{81}(F_2-I)=0\pmod{81}$.
It defines an actual local coefficient quotient
$M\twoheadrightarrow\mathbb Z/81$ with trivial
Frobenius action, and reads the coordinate $x_{12}$.
The script records its entire row and checks both
primitivity and the annihilation identity directly.
The surviving line changes this coordinate by
$27$ or $54$. Reduction of the local coinvariant
group modulo $3^r$, $r\le3$, kills that line; reduction
modulo $81$ retains it. This is a statement about
local cohomology at two, not about all global finite
coefficient reductions or the order of an entire
finite geometric quotient.
The additional datum is the actual Frobenius value under
(B.38.15), in its fixed normalization. B.39 supplies an
intrinsic class-two quotient detecting it, and B.40
evaluates it through the actual geometric extension of
the multiplication-three elliptic cover. Thus the full
refinement is selected without relying on an unevaluated
point count or on the inherited quotient.

By contrast, B.32 supplies representatives of $u,v$
inside the ambient tripod commutator group $F'$.
Every bracket generator of $\Lambda^2V$ has a commutator
representative in $N'\subset F'$.
Therefore each generator $[h,[n,n']]$ of $\mathcal J$
belongs to $F''$.
Changing a representative only adds a term in
$\gamma_4N$, so this proves the exact coefficient
vanishing
\[
 \mathcal J\subset\ker\left(
 M\longrightarrow
 \gamma_3\bigl(N/(F''\gamma_4N)\bigr)(-1)\right).             \tag{B.38.16}
\]
The finite quotient in B.22 is obtained inside this
inherited metabelian ambient group. It also kills
$\mathcal J$, and consequently kills every member
of the surviving homogeneous line.
The three formal integral refinements therefore
have identical images at every Frobenius in that
prescribed finite quotient. This strengthens the
failure of that proposed test under the explicit
geometric constraints considered here, while
leaving the actual choice among the three formal
classes open.
\end{proof}
